\chapter{Conclusions}
\label{ch:conclusions}

\section{Summary of the Work}
\label{sec:conclusions_summary}

Chapter~\ref{ch:introduction} established three foundational research questions under a strict constraint: an LLM performing semantic routing on natural-language workforce requests while excluding personal data from its active context, yet maintaining the operational precision required for an administrator to act upon the output without manual verification. The subsequent chapters addressed these inquiries empirically through the implementation and validation of a functional system rather than relying exclusively on theoretical arguments.

The first question investigated the distribution of responsibilities among the client, the anonymisation layer, the LLM, and the business logic server to structurally minimise PII exposure rather than depending on operational policy. The four-tier topology introduced in Chapter~\ref{ch:architecture} satisfies this requirement. The Model Context Protocol (MCP) client anonymises requests via the Anonymisation Service prior to external transmission. Specifically, GLiNER detects sensitive entities, the Anonymisation Service substitutes them with typed and numbered tokens, and the mappings are persisted within Redis under unique context identifiers. Consequently, \texttt{ask\_antares} receives solely the tokenised representation and the context identifier. \texttt{JobTools} restores original values server-side at the point of execution after Claude generates its routing decision (Section~\ref{sec:arch_privacy_by_design}).

By design, no component upstream of \texttt{JobTools} possesses both a token and its underlying plaintext value, ensuring the model interacts exclusively with tokenised artifacts. Chapter~\ref{ch:usecases} evaluated the practical effectiveness of this architecture, revealing that detection recall governs overall privacy guarantees. Scanning fifty test prompts for literal plaintext strings exposed unmasked PII in two instances. Both cases stemmed from specific detection failures by GLiNER---wherein entities were overlooked entirely rather than leaking through already-established tokens (Section~\ref{subsec:usecases_llm_audit}).

The second question examined the detection tier, evaluating whether a zero-shot, multilingual Named Entity Recognition (NER) model could reach practically usable recall on Italian business terminology after the rule-based alternative, Presidio, proved inadequate (Section~\ref{subsec:impl_gliner_pivot}). Across 82 labelled entity spans distributed through 50 prompts, GLiNER attained 93.8\% precision and 91.5\% recall on the coverage metric, which dictates whether PII can infiltrate the LLM context (Table~\ref{tab:gliner_metrics}). The seven missed spans were not randomly distributed; four shared a common structural origin linked to Italian date-range expressions structured as ``dal X al Y'', as detailed in Table~\ref{tab:gliner_fn}.

The third question analysed the operational cost that reversible tokenisation imposes on system correctness. Twelve end-to-end test scenarios across three distinct use cases---each verified against an independent reference implementation sharing no common codebase---exhibited complete parity (Tables~\ref{tab:uca_e2e}, \ref{tab:ucb_e2e}, and \ref{tab:ucc_e2e}). Routing accuracy over tokenised inputs reached 95.2\% using the log-diffing evaluation methodology. Manual inspection of the single discrepancy confirmed that Claude selected the appropriate tool before encountering a legitimate business logic exception prior to log generation, indicating a methodological limitation rather than a routing failure (Table~\ref{tab:routing_summary}). Furthermore, running 21 distinct prompts across five repetitions each at a sampling temperature of 1.0 yielded identical routing selections in every trial. Within this test set, the privacy layer did not cost measurable correctness: routing and end-to-end results matched the reference implementation and stayed stable across repeated sampling, though the test set is 21 hand-built prompts, not a claim covering the full space of real user phrasing.

The validation phase also uncovered minor discrepancies in earlier architectural descriptions. Chapter~\ref{ch:architecture} originally stated that GLiNER split the canonical Bergamo prompt into two location tokens; however, telemetry data for Table~\ref{tab:routing_summary} confirmed that three tokens were generated due to compound site phrasing. While minor, this correction underscores a broader principle: architectural specifications that remain unchecked against running systems frequently harbor inaccuracies discoverable solely through direct runtime observation.

These results warrant careful qualification. GLiNER's coverage recall of 91.5\% dictates the practical ceiling of the system's privacy guarantees, as undetected entities reach the model unmasked (Section~\ref{subsec:usecases_llm_audit}). Furthermore, the Redis context store lacks an active time-to-live configuration, and server-side deanonymisation fails to purge read mappings, causing token associations to persist throughout the process lifecycle beyond failed or abandoned requests (Section~\ref{sec:discussion_security_complexity}). Additionally, the absence workflow for Use Case C remained unexercised due to test windows falling outside the mock dataset's absence-ticket interval (Section~\ref{sec:usecases_holiday_request}). Finally, evaluations utilised a mock Miorelli backend rather than a live production tenant, a scoping decision agreed upon with the enterprise partner to prioritise privacy-layer validation (Section~\ref{sec:intro_contributions}).

\section{Final Remarks on Enterprise AI Adoption}
\label{sec:conclusions_final_remarks}

The General Data Protection Regulation (GDPR) draws a definitive legal distinction between pseudonymisation and anonymisation. As outlined in Section~\ref{subsec:background_pii_definitions}, substituting a proper name with \texttt{<PERSON\_0001>} relocates personal data rather than destroying it---transferring information from an uncontrolled third-party provider infrastructure to an internal Redis repository managed by the organisation~\cite{gdpr2016}. This distinction dictates how the conclusions of this work generalise. Mitigating LLM exposure does not extinguish core regulatory obligations; instead, it shifts accountability to the storage layer of the token mappings, which demands identical engineering rigor, access control, and retention management as raw personal data~\cite{enisa2019pseudonymisation}. The absence of an automated time-to-live policy within the project's Redis store is exactly this risk in practice: a boundary between models and sensitive data does not, on its own, secure the data store behind that boundary.

The Model Context Protocol's separation of host, server, and model roles~\cite{anthropic2024mcp} provided an established framework for integrating this boundary without building infrastructure from scratch. The empirical contribution of this thesis lies in validating that these enforcement mechanisms hold under operational testing rather than relying solely on protocol documentation (Section~\ref{sec:usecases_validation}). Systematic reviews of the broader MCP ecosystem indicate that few production deployments enforce these available safeguards~\cite{hou2025mcpsecurity}. Antares AI successfully implements model isolation for PII while sharing common operational vulnerabilities---such as local log retention, unencrypted storage at rest, and insufficient access controls---identified across the wider MCP landscape.

Section~\ref{sec:background_sota} did not find a system in the literature it reviewed that combines local detection, reversible tokenisation, external semantic routing, and deterministic server-side execution alongside a validated correctness benchmark. Antares AI closes part of that gap, within the scope of that same review. The validation dataset reflects one specific enterprise platform rather than a cross-domain general-purpose benchmark. However, the foundational methodology remains broadly reusable: detecting PII locally, tokenising reversibly, routing via tokens externally, executing within trusted boundaries, and directly measuring detection recall. The specific numerical accuracy metrics are tied to this dataset and prompt distribution, whereas the overarching architectural pattern offers a viable blueprint for alternative domains and backends.

The replacement of Presidio with GLiNER (Section~\ref{subsec:impl_gliner_pivot}) carries implications extending beyond technical performance. The initial architecture was approved using Presidio, yet subsequent testing against real Italian enterprise inputs necessitated its replacement due to insufficient recall. Implementing privacy-by-design~\cite{cavoukian2009privacy} in practice required discarding an approved component following empirical validation. This mirrors counterfactual analyses of privacy failures at major technology firms, where abstract commitments frequently fail when weighed against specific engineering decisions~\cite{rubinstein2013counterfactual}. Similarly, Use Case C's authorisation gate depends on hardcoded identities that must be re-engineered once real identity providers are introduced (Section~\ref{sec:discussion_production_migration}). These observations suggest that enterprise privacy engineering is less constrained by initial architectural planning than by the organisational willingness to revise components when confronted with operational realities.

The findings support one precise claim: privacy and capability are compatible in language-model-powered enterprise systems under three preconditions. The boundary separating models from sensitive data must be structural rather than presumptive. Detection recall must be measured, not assumed. Downstream business logic must stay deterministic and auditable rather than delegated to a probabilistic model. Antares AI met these three conditions within one workforce management context, on the dataset and prompts this thesis tested. Whether the same conditions hold in other domains, languages, or with other model providers is untested and open for future work.